% =========================================================
% CHAPITRE 1 — DESCRIPTION DU PROJET
% =========================================================

\section{Introduction}

Le développement durable et la réduction des déchets plastiques constituent
aujourd'hui des enjeux importants sur les plans environnemental, économique
et technologique. Parmi les matières plastiques les plus utilisées, le
polyéthylène téréphtalate (PET) occupe une place importante, notamment dans
la fabrication des bouteilles et des emballages.

Le PET présente cependant l'avantage d'être recyclable et peut être valorisé
afin de produire de nouveaux matériaux. Dans le domaine de la fabrication
additive, le plastique recyclé peut notamment être transformé en filament
destiné à l'impression 3D. Cette approche permet de donner une seconde vie
aux déchets plastiques tout en réduisant la consommation de matière première.

C'est dans ce contexte que s'inscrit le projet \textit{3awedlou}, une
solution IoT dédiée au recyclage du PET et à sa transformation en filament
pour impression 3D. Le projet associe une partie matérielle permettant les
différentes étapes du processus de transformation à une architecture
logicielle destinée à la supervision et au suivi des données de la machine.

L'objectif de ce projet ne se limite donc pas à la réalisation d'une
machine de transformation du plastique. Il vise également à intégrer les
principes de l'Internet des Objets afin de rendre le système connecté,
supervisable et capable de fournir des informations sur son fonctionnement.

\section{Problématique}

La gestion des déchets plastiques représente un défi important en raison de
leur volume, de leur durée de dégradation et de leur utilisation massive
dans de nombreux secteurs. Le PET, largement présent dans les bouteilles
plastiques, constitue une matière particulièrement intéressante pour la
valorisation en raison de ses propriétés et de sa recyclabilité.

Cependant, transformer des déchets plastiques en un matériau utilisable
pour l'impression 3D nécessite plusieurs opérations successives. Le
plastique doit notamment être préparé, chauffé, fondu et extrudé de manière
contrôlée afin d'obtenir un filament présentant des caractéristiques
adaptées à son utilisation.

Une difficulté supplémentaire réside dans le contrôle et la supervision de
ces différentes opérations. La température, le fonctionnement du moteur,
l'état de la machine et les paramètres de production doivent pouvoir être
surveillés afin d'améliorer la fiabilité du système et d'identifier
rapidement les éventuels problèmes.

La problématique principale de ce projet peut ainsi être formulée comme
suit :

\begin{quote}
\textit{Comment concevoir et réaliser une machine connectée capable de
transformer des déchets PET en filament pour impression 3D tout en
permettant le contrôle, la supervision et la collecte des données de son
fonctionnement ?}
\end{quote}

Cette problématique implique donc à la fois des considérations mécaniques,
électroniques, informatiques et liées à l'Internet des Objets.

\section{Solution Proposée}

Afin de répondre à cette problématique, le projet propose la conception et
la réalisation d'un système de recyclage du PET intégrant une architecture
IoT.

La solution repose sur une machine physique chargée d'assurer la
transformation du plastique. La partie matérielle comprend notamment un
système de chauffage et de contrôle de température, un moteur pas à pas
destiné à l'extrusion du matériau ainsi que différents éléments
électroniques permettant le pilotage et la mesure des paramètres de
fonctionnement.

La commande du système est assurée par un microcontrôleur Arduino Mega.
Celui-ci prend en charge une partie du contrôle local de la machine,
notamment la gestion du moteur et du système de chauffage. Un ESP32 est
également intégré à l'architecture afin d'assurer les fonctions de
communication et de connectivité du système.

La communication entre les différents éléments permet de séparer les
fonctions de contrôle temps réel de la machine et les fonctions de
communication avec les services logiciels. Cette architecture facilite
ainsi l'évolution du système et son intégration dans une infrastructure
IoT.

Sur le plan logiciel, le système est associé à une application permettant
de superviser la machine et d'exploiter les données collectées. Le backend
constitue l'intermédiaire entre les équipements connectés et les
applications clientes. Il permet notamment de gérer les utilisateurs, les
machines, les sessions de fonctionnement, les données de télémétrie et les
informations relatives à la production.

La communication IoT est basée sur le protocole MQTT, particulièrement
adapté aux systèmes connectés grâce à son modèle de communication
publish/subscribe et à sa faible surcharge. Les données provenant de la
machine peuvent ainsi être transmises au backend afin d'être enregistrées
et exploitées par les interfaces de supervision.

L'architecture globale du système peut être représentée selon plusieurs
niveaux : la machine et ses capteurs/actionneurs constituent la couche
physique, les microcontrôleurs assurent le contrôle et la communication,
le backend assure la gestion des données et les applications web et
mobiles permettent leur visualisation et leur exploitation.

Cette approche permet de transformer une machine de recyclage classique en
un système connecté capable de fournir des informations en temps réel sur
son fonctionnement. Elle constitue également une base évolutive pour
l'ajout de nouvelles fonctionnalités telles que l'analyse des données,
l'amélioration du contrôle du processus ou encore la maintenance
préventive.

\section{Conclusion}

Dans ce chapitre, nous avons présenté le contexte général du projet
\textit{3awedlou}, ainsi que la problématique liée à la valorisation des
déchets PET et à leur transformation en filament pour impression 3D.

La solution proposée repose sur l'association d'une machine de recyclage,
de systèmes embarqués et d'une architecture IoT permettant le contrôle, la
communication et la supervision du système.

Après avoir présenté le projet et ses objectifs généraux, le chapitre
suivant sera consacré à l'étude détaillée des besoins ainsi qu'à la
conception matérielle et logicielle de la solution.